\begin{tabular}{lrrrrr}
\hline\hline
& 1980 & 1990 & 2000 & 2010 & 2020 \\
\hline
Municipalities in the delineation & 1,656 & 1,656 & 1,656 & 1,656 & 1,655 \\
Commuting zones & 385 & 307 & 261 & 231 & 223 \\
Single-municipality zones & 72 & 35 & 16 & 9 & 10 \\
Municipalities in the largest zone & 31 & 29 & 30 & 27 & 29 \\
Non-contiguous zones & 0 & 0 & 0 & 0 & 0 \\[3pt]
Smallest zone, resident labour force & 286 & 788 & 523 & 530 & 475 \\
Average zone, resident labour force & 120,006 & 168,711 & 202,706 & 202,694 & 207,532 \\
Largest zone, resident labour force & 4,761,443 & 5,927,239 & 6,000,065 & 4,929,463 & 5,624,897 \\[3pt]
Smallest zone area (sq.\,km) & 16 & 24 & 109 & 184 & 100 \\
Average zone area (sq.\,km) & 944 & 1,184 & 1,393 & 1,574 & 1,630 \\
Largest zone area (sq.\,km) & 3,464 & 4,429 & 4,429 & 5,744 & 5,744 \\[3pt]
Compactness of zones & 0.33 & 0.32 & 0.31 & 0.30 & 0.30 \\
Compactness of municipalities & 0.41 & 0.41 & 0.41 & 0.41 & 0.41 \\[3pt]
United States, compactness of commuting zones & \multicolumn{5}{c}{0.40} \\
United States, compactness of counties & \multicolumn{5}{c}{0.53} \\
\hline\hline
\end{tabular}
